\section{Data and measurement}
\subsection{Survey}
The recorded survey interviews span 9 December 2025 to 17 February 2026. Fundación Capital recruited migrant women leaders who administered the survey, approaching women they knew and following referrals from respondents. Women who had not participated in Fundación Capital programs could take part. Recruitment was intended to reach adult women with international migration experience, including Colombian return migrants, with approximately three months' residence after arrival or return. The programmed screen used consent, gender, age and whether the respondent or someone in her household had received a remittance in the previous 12 months. The instrument did not test every intended recruitment criterion. The screen did not establish personal receipt. A separate question asked whether the respondent was the household's principal recipient or manager.

The analysis retains the historical 423-record cohort and all original manual exclusions, including test and other non-official submissions. The documented coded filters take 490 records to 423. That accounting is not a response rate, and the record does not establish where each manual exclusion occurred relative to that export. Referral links, network degrees and household links were not systematically recorded. We therefore treat the survey as a nonprobability referral sample, without claiming respondent-driven sampling weights or population estimates \citep{gileHandcock2010}. We cannot verify 423 distinct households. Recorded residence is not the interview site; some residence codes changed meaning or availability across form versions.

The first outcome summarizes what respondents say happens to incoming remittance proceeds in general (q6\_13): any reported digital retention or use versus all cash. The question was asked on the eligible route whether or not the respondent said she was the principal recipient or manager (q6). Its answers describe reported disposition in the household-screened sample, not disposition among verified personal recipients. An all-cash answer does not imply that the money was first received digitally and then withdrawn. The five choices are also retained in questionnaire order. Two choices describe overlapping proportions, so we neither reconstruct shares of pesos nor rank the five responses as a scale. The latest receipt channel and 12-month receipt frequency are separate items. A transfer received through a wallet may be withdrawn; a transfer received in cash need not imply that the recipient never uses domestic digital payments. The frequency item contains a combined response for nonapplicability or not remembering, which does not establish nonreceipt.

The second outcome records whether respondents made or received a domestic digital payment or transfer in the previous 60 days. The denominator is valid yes/no answers. Refusal and other unscored states are not No. The follow-up item uses five approximate operation bins among affirmative users. Its scorable-bin base excludes uncertainty; the response-state denominator includes all affirmative users. Overlapping endpoints do not yield exact counts. The fee-visibility item concerns the last operation in a 90-day window but is asked only on the affirmative 60-day branch. It does not describe all 90-day users. Neither main outcome traces a payment from sender to final use.

The models use three fixed summaries of survey answers. The phone, internet and data score (IAT3) averages habitual phone use, internet-use frequency and whether running out of data or balance prevented the respondent from using financial services or communicating in the last month. The last item was asked only of regular smartphone or basic-phone users. The score does not describe someone who never reached that question, and part of it is close to the behavior the models describe. The ability and confidence score (IADT2) averages reported ability to send or receive money by phone and confidence in avoiding online fraud or suspicious messages. The transfer item was also asked only of regular phone users; neither item tests actual skill. The two-condition onboarding score (IAFF2) averages whether the phone number normally used for a digital account or wallet was registered in the respondent's name and whether she reported address proof in her own name that a financial institution would accept. Neither question verifies an account, a provider's decision or effective control. Each score requires a scorable answer for every component. A separately labeled IADT3 alternative adds reported app-installation ability, asked only on the smartphone route. The supplement defines the scores and distinguishes bases for possible comparisons; it does not report new correlation results. Correlations among these measures would not establish construct validity.

The descriptive denominators follow what was asked and answered. A question may have been unavailable in the recorded form, not applicable because of a skip, applicable but unanswered, or answered with an explicit uncertainty or refusal option. Those states have different meanings. The version-aware accounting preserves them rather than assigning zeros, borrowing response codes from another form or assuming that complete cases are selected at random \citep{rubin1976}. We distinguish the admitted cohort, item-specific denominators and model samples.

Table~\ref{tab:academic-sample} summarizes the age, recorded residence and reported education of the recruited sample. Table~\ref{tab:academic-measurement} shows the valid denominators and shares for the two outcomes, and the scorable counts, means and standard deviations for the three regression measures. The scores are fixed summaries of observed answers, not estimates of a common latent trait. A higher mean on one score than on another does not show greater inclusion in that domain.

\input{"\CFIDescriptivesDir sample_characteristics.tex"}
\input{"\CFIDescriptivesDir measurement_summary.tex"}

\subsection{Interviews, focus groups and documentary comparison}

The fieldwork archive records 33 semistructured interviews with Venezuelan migrant women across five sites (Bogotá, Soacha, Cali, Medellín and Villavicencio), including three in Villavicencio. The retained thematic synthesis reports a base of 30 interviews. The available coding documentation does not establish whether the three Villavicencio interviews fed that synthesis. The qualitative strand also records six focus groups involving 47 participants; the interview and group counts do not establish the number of distinct people across activities.

The earlier thematic analysis combined the study's conceptual domains with themes developed through repeated reading of interviews and focus-group material. This paper re-examines Angélica Cardozo Cadavid's written synthesis and memo. We distinguish a reported experience from a question or exercise in the field guides, then compare accounts by direction, segment of the transaction and the roles of the people involved. The original discussions and interviews are the source material behind the synthesis; the current paper does not recode transcripts, count mentions or link qualitative cases to survey records. The technical supplement's \emph{Provenance of the qualitative evidence} identifies the source sections for the claims we draw on.

The original paper also draws on Colombian institutional and regional expert interviews and reviews regulatory, central-bank and policy documents. The expert evidence helps identify reported implementation problems; official sources document selected rules and system designs. Neither establishes that a particular respondent could open an account, used a specific rail or resolved a complaint. The regional comparison examines counterpart systems, not comparable groups of women in five countries.
